\section*{Capítulo \ref{ch:b2:pt-paths} --- Termobarometría y trayectorias P--T--t}

\begin{solution}{exo:b2:pt-paths:1}
(a) A zona de la clorita, B zona de la biotita, C zona del granate, D zona de la estaurolita, E zona de la cianita,
F zona de la sillimanita. (b) Cinco: las isogradas de la biotita, el granate, la estaurolita, la cianita y la
sillimanita, cada una entre dos muestras sucesivas, con el grado aumentando de A a F. (c) Una isograda
solo dice que una reacción tuvo lugar; un gradiente de campo necesita presiones y temperaturas, que
proceden de reacciones calibradas o de la termobarometría de la asociación de pico de cada
roca.
\end{solution}

\begin{solution}{exo:b2:pt-paths:2}
(a) $(X_{\text{Mg}}/X_{\text{Fe}})_{\text{grt}} = 0.10/0.80 = 0.125$,
$(X_{\text{Mg}}/X_{\text{Fe}})_{\text{bt}} = 0.45/0.50 = 0.900$; $K_D = 0.1389$, $\ln K_D =
-1.9741$. (b) $3R = 24.942$; el denominador es $20.19 + 24.942\times1.9741 = 69.428$; a
\qty{5}{kbar}, $T = 54185/69.428 = \qty{780.5}{K} = \qty{507}{\celsius}$; a
\qty{7}{kbar}, $T = (52630 + 2177)/69.428 = \qty{789.4}{K} = \qty{516}{\celsius}$.
(c) Dos kilobares la cambian solo \qty{9}{K}: un termómetro.
\end{solution}

\begin{solution}{exo:b2:pt-paths:3}
(a) $z = P/(\bar\rho g) = 6.0\times10^8/(2800\times9.81) = \qty{21.8}{km}$.
(b) $(580 - 10)/21.8 = \qty{26}{\celsius\per\km}$. (c) A \qty{580}{\celsius} la línea
está en $3.67 + 0.8\times2.10 = \qty{5.35}{kbar}$; a \qty{6.0}{kbar} la roca está del
lado de alta presión: cianita.
\end{solution}

\begin{solution}{exo:b2:pt-paths:4}
(a) Reactivos $2\times140.693 + 3\times18.015 = \qty{335.431}{g}$; productos $277.112 +
58.320 = \qty{335.432}{g}$: ajustada salvo redondeo. (b) $2\times43.65 =
\qty{87.30}{cm^3}$ se convierten en $107.50 + 24.63 = \qty{132.13}{cm^3}$: $+\qty{51.4}{\percent}$.
(c) Olivino $281.386/87.30 = \qty{3.223}{g.cm^{-3}}$; productos $335.432/132.13 =
\qty{2.539}{g.cm^{-3}}$. (d) La roca aumenta su volumen en la mitad y se vuelve más ligera; además
almacena agua en sus minerales.
\end{solution}

\begin{solution}{exo:b2:pt-paths:5}
(a) Medias $-1.206$ a \qty{1072.15}{K} y $-1.782$ a \qty{823.15}{K}. Con $1/T =
9.3271\times10^{-4}$ y $1.21485\times10^{-3}$, $a = (-1.206 + 1.782)/(9.3271\times10^{-4}
- 1.21485\times10^{-3}) = -2041.5$ K y $b = -1.206 + 2041.5/1072.15 = 0.6982$.
(b) $\Delta_r H + 2070\times0.311 = 2041.5\times24.942 = \qty{50920}{J}$, luego $\Delta_r H =
\qty{50276}{J}$; $\Delta_r S = 0.6982\times24.942 = \qty{17.41}{J.K^{-1}}$. (c) A
\qty{5}{kbar}: para $K_D = 0.18$, 588.1 frente a \qty{587.5}{\celsius}; para $0.10$, 419
frente a \qty{425}{\celsius}; para $0.30$, 819 frente a \qty{806}{\celsius}. Dentro del intervalo
de los experimentos las dos calibraciones coinciden, aunque sus constantes difieren en
\qty{2.4}{kJ} y \qty{2.8}{J.K^{-1}}; los errores de $\Delta_r H$ y de $\Delta_r S$
se compensan allí y crecen fuera de él. Una calibración no debe usarse lejos
de las condiciones de sus experimentos.
\end{solution}

\begin{solution}{exo:b2:pt-paths:6}
(a) Intercambio: $(113.12 + 154.30) - (115.32 + 149.65) = \qty{+2.45}{cm^3} =
\qty{0.245}{J.bar^{-1}}$, frente a 0.311. GASP: $125.28 + 2\times44.15 + 22.69 -
3\times100.79 = \qty{-66.10}{cm^3} = \qty{-6.61}{J.bar^{-1}}$, frente a $-6.21$, el valor
de la calibración a \qty{800}{\celsius} y \qty{10}{kbar}, donde la dilatación térmica y la
compresión cambian los volúmenes. En ambos casos la variación de GASP es de veinte a treinta veces
la del intercambio. (b) $\partial T/\partial P = 0.311/62.96 = \qty{4.9e-3}{K/bar}$:
\qty{4.9}{K} por kilobar. (c) La pendiente de la línea es $22.80 + 8.314\times3\ln 0.20/6.21
= 22.80 - 6.46 = \qty{16.3}{bar/K}$; \qty{17}{K} desplazan la presión
\qty{280}{bar}.
\end{solution}

\begin{solution}{exo:b2:pt-paths:7}
(a) $P_0 = 22.80\times600 - 1093 = \qty{12587}{bar}$; $3\ln(0.08/0.30) = -3.9653$; $RT =
8.314\times873.15 = \qty{7259.4}{J}$; $P = 12587 + 7259.4\times(-3.9653)/6.21 = 12587 -
4635 = \qty{7952}{bar}$. (b) $3\ln(0.08/0.35) = -4.4277$, $P = 12587 - 5176 =
\qty{7411}{bar}$: \qty{0.54}{kbar} menos, de modo que la plagioclasa debe analizarse con cuidado.
(c) $7952\times10^5/(2800\times9.81) = \qty{29.0}{km}$.
\end{solution}

\begin{solution}{exo:b2:pt-paths:8}
(a) $z = P/(\bar\rho g)$: \qty{118.9}{km}, \qty{34.0}{km} y \qty{15.3}{km}.
(b) $(118.9 - 34.0)/2.2 = \qty{38.6}{km/Myr} = \qty{3.9}{cm.yr^{-1}}$ y $(34.0 -
15.3)/1.1 = \qty{17.0}{km/Myr} = \qty{1.7}{cm.yr^{-1}}$, algo por encima de los valores publicados,
3.4 y 1.6: su conversión de presión en profundidad da diferencias de profundidad algo
menores. (c) $\sigma = \sqrt{0.9^2 + 0.9^2} = \qty{1.3}{Myr}$, luego $2.2 \pm 1.3$
millones de años y una velocidad entre $84.9/3.47 = 24$ y $84.9/0.93 = \qty{91}{km/Myr}$,
de 2.4 a \qty{9.1}{cm.yr^{-1}}: rápida en cualquier caso. (d) Incluso la velocidad más baja es unas cincuenta veces
los \qty{0.5}{mm.yr^{-1}} del modelo: la erosión sola no puede lograrlo.
\end{solution}

\begin{solution}{exo:b2:pt-paths:9}
(a) $\partial T/\partial\ln K_D = 3RT/(\Delta_r S - 3R\ln K_D) = 24.942\times860.6/62.96 =
\qty{341}{K}$; $\partial T/\partial P = 0.311/62.96 = \qty{4.9}{K}$ por kilobar.
(b) $\sigma_T = \sqrt{(341\times0.05)^2 + 4.9^2} = \sqrt{17.0^2 + 4.9^2} =
\qty{17.7}{K}$. (c) Domina el error analítico de $K_D$; por sí solo supone dos tercios de
los \qty{25}{\celsius} de las mejores calibraciones, que incluyen además los errores de la
propia calibración.
\end{solution}

\begin{solution}{exo:b2:pt-paths:10}
(a) Granate $0.10/0.75 = 0.1333$, biotita $0.42/0.52 = 0.8077$, $K_D = 0.1650$, $\ln K_D =
-1.8018$, denominador $20.19 + 24.942\times1.8018 = 65.130$; $3\ln(0.07/0.25) = -3.8189$.
A \qty{6000}{bar}: $T = 54496/65.130 = \qty{836.7}{K} = \qty{563.6}{\celsius}$, después $P
= 22.80\times563.6 - 1093 + 8.314\times836.7\times(-3.8189)/6.21 = \qty{7479}{bar}$.
A \qty{7479}{bar}: $T = \qty{570.6}{\celsius}$, $P = \qty{7603}{bar}$. A
\qty{7603}{bar}: $T = \qty{571.2}{\celsius}$, $P = \qty{7614}{bar}$; una vez más, $T =
\qty{571.3}{\celsius}$. La roca se equilibró a unos \qty{571}{\celsius} y
\qty{7.6}{kbar}. (b) A \qty{571}{\celsius} la línea cianita $=$ sillimanita está en $3.67
+ 0.71\times2.10 = \qty{5.16}{kbar}$: cianita. (c) El barómetro se calibró para la
cianita; con sillimanita la reacción, y su línea, son otras, ya que los dos polimorfos no
tienen la misma energía de Gibbs.
\end{solution}

\begin{solution}{exo:b2:pt-paths:11}
(a) $\tau = L^2/(\pi^2\kappa)$: para \qty{35}{km}, $1.225\times10^9/(9.870\times10^{-6}) =
\qty{1.24e14}{s} = \qty{3.9}{Myr}$; para \qty{70}{km}, cuatro veces más,
\qty{15.7}{Myr}. (b) Setenta millones de años son de cuatro a dieciocho veces esos tiempos: el calor
atraviesa la corteza engrosada mucho más deprisa de lo que la erosión la retira. (c) La roca
alcanza su máxima profundidad de golpe, en el apilamiento; el calor llega a lo largo de las
decenas de millones de años siguientes, mientras la erosión la va subiendo despacio: la temperatura sigue
subiendo un tiempo mientras se produce la descompresión. (d) Una erosión en siete millones de años
subiría la roca antes de que la conducción la hubiera calentado: un pico térmico más bajo, alcanzado
pronto, a mayor presión, y una trayectoria fría de vuelta a la superficie.
\end{solution}

\begin{solution}{exo:b2:pt-paths:12}
(a) La corteza pesa $35000\times2800\times9.81 = \qty{0.961}{GPa}$; los
\qty{1.839}{GPa} restantes requieren $1.839\times10^9/(3300\times9.81) = \qty{56.8}{km}$ de manto: profundidad
\qty{91.8}{km}, como los \qty{90}{km} o más del estudio original. (b) $(4.0 -
0.961)\times10^9/(3300\times9.81) = \qty{93.9}{km}$, profundidad \qty{128.9}{km}.
(c) $700/91.8 = \qty{7.6}{\celsius\per\km}$: una razón muy alta entre presión y
temperatura, la serie de las zonas de subducción. (d) $(91.8 - 35)/34 = \qty{1.7}{Myr}$.
\end{solution}

\begin{solution}{pb:b2:pt-paths:1}
\noindent\textbf{Parte I.}

\textbf{1.} Núcleo del granate $0.080/0.760 = 0.1053$, biotita $0.40/0.58 = 0.6897$; $K_D =
0.1527$, $\ln K_D = -1.8793$.

\textbf{2.} Borde $0.120/0.720 = 0.1667$, biotita de la matriz $0.46/0.53 = 0.8679$; $K_D =
0.1921$, $\ln K_D = -1.6497$.

\textbf{3.} $3R = 24.942$; numerador $52630 + 0.311\times6000 = 54496$; denominador $20.19
+ 24.942\times1.8793 = 67.064$; $T = \qty{812.6}{K} = \qty{539.5}{\celsius}$.

\textbf{4.} Denominador $20.19 + 24.942\times1.6497 = 61.337$; $T = 54496/61.337 =
\qty{888.5}{K} = \qty{615.3}{\celsius}$.

\textbf{5.} $0.311\times1000/67.064 = \qty{4.6}{K}$ para el núcleo y $0.311\times1000/61.337
= \qty{5.1}{K}$ para el borde.

\textbf{6.} Un termómetro se aplica a minerales que estuvieron juntos en equilibrio. La
biotita incluida en el núcleo quedó aislada del resto de la roca cuando creció el núcleo;
la biotita de la matriz intercambió por última vez con el borde. Emparejar el núcleo con la matriz mezclaría
dos momentos.

\noindent\textbf{Parte II.}

\textbf{7.} Núcleo $22.80\times539.5 - 1093 = \qty{11208}{bar}$; borde $22.80\times615.3 -
1093 = \qty{12936}{bar}$.

\textbf{8.} $3\ln(0.050/0.25) = -4.8283$; $RT = 8.314\times812.65 = \qty{6756.4}{J}$; $P =
11208 - 6756.4\times4.8283/6.21 = 11208 - 5253 = \qty{5955}{bar}$.

\textbf{9.} $3\ln(0.060/0.28) = -4.6213$; $RT = 8.314\times888.45 = \qty{7386.6}{J}$; $P
= 12936 - 5497 = \qty{7439}{bar}$.

\textbf{10.} A \qty{7439}{bar}: $T = (52630 + 2314)/61.337 = \qty{895.8}{K} =
\qty{622.6}{\celsius}$; a esa temperatura $P = 13102 - 5542 = \qty{7560}{bar}$; después $T
= \qty{623.2}{\celsius}$. El borde creció a unos \qty{623}{\celsius} y \qty{7.56}{kbar}.

\textbf{11.} A \qty{5954}{bar}, $T = \qty{539.2}{\celsius}$ y $P = \qty{5950}{bar}$: el
núcleo creció a unos \qty{539}{\celsius} y \qty{5.95}{kbar}. La primera estimación,
\qty{6}{kbar}, ya era cercana, y el termómetro solo se desplaza \qty{4.6}{K} por kilobar.

\textbf{12.} Núcleo $5950\times10^5/(2800\times9.81) = \qty{21.7}{km}$; borde
$7560\times10^5/27468 = \qty{27.5}{km}$.

\textbf{13.} A \qty{539}{\celsius} la línea está en $3.67 + 0.39\times2.10 =
\qty{4.49}{kbar}$, por debajo de \qty{5.95}{kbar}: cianita, coherente con la cianita incluida
en el núcleo. A \qty{623}{\celsius} está en $5.77 + 0.23\times2.09 = \qty{6.25}{kbar}$,
por debajo de \qty{7.56}{kbar}: de nuevo cianita.

\noindent\textbf{Parte III.}

\textbf{14.} $+\qty{84}{\celsius}$ y $+\qty{1.6}{kbar}$ (unos \qty{6}{km} más abajo): la
roca se estaba enterrando y calentando, en su \omterm{def:b2:pt-paths:paths}{trayectoria prógrada}.

\textbf{15.} A \qty{623}{\celsius} o más, por debajo de \qty{6.25}{kbar}: después de crecer el borde la
roca perdió al menos \qty{1.3}{kbar}, unos \qty{5}{km} de recubrimiento, sin dejar de estar
al menos igual de caliente. Se exhumó mientras aún estaba caliente.

\textbf{16.} La clorita es un mineral hidratado de bajo grado: registra la parte final de la \omterm{def:b2:pt-paths:paths}{trayectoria
retrógrada}, el enfriamiento a poca profundidad, con agua.

\textbf{17.} Enterramiento y calentamiento hasta el borde, descompresión sin enfriamiento, y después enfriamiento:
dibujada con la temperatura hacia la derecha y la presión hacia arriba, la trayectoria gira en sentido horario. El
pico de presión llegó primero, en el borde; el pico térmico a la vez o más tarde, durante la
descompresión.

\textbf{18.} Sí: el modelo también alcanza primero su pico de presión y su pico térmico,
\qty{593}{\celsius} a \qty{7.8}{kbar}, durante el ascenso, cerca de las condiciones del borde
del esquisto. Un engrosamiento seguido de una erosión lenta explica la trayectoria del esquisto.

\noindent\textbf{Parte IV.}

\textbf{19.} $(623 - 405)/(412.0 - 398.0) = 218/14.0 = \qty{15.6}{\celsius}$ por millón
de años.

\textbf{20.} La roca se enfrió alrededor de una vez y media más deprisa de lo supuesto. Un mineral
que se enfría más deprisa tiene menos tiempo para perder su argón, de modo que se cierra a una temperatura más alta:
la verdadera temperatura de cierre está algo por encima de \qty{405}{\celsius}, y el cierre
algo más abajo.

\textbf{21.} $(405 - 10)/30 = \qty{13.2}{km}$.

\textbf{22.} $27.5 - 13.2 = \qty{14.3}{km}$ (con las profundidades sin redondear, 27.52 y 13.17,
\qty{14.35}{km}).

\textbf{23.} $\sqrt{1.5^2 + 1.0^2} = \qty{1.8}{Myr}$: $14.0 \pm \qty{1.8}{Myr}$, luego una velocidad
entre $14.35/15.8 = 0.91$ y $14.35/12.2 = \qty{1.18}{km/Myr}$.

\textbf{24.} $14.35/14.0 = \qty{1.03}{km}$ por millón de años, es decir,
\textbf{\qty[text-series-to-math=true]{1.0}{mm.yr^{-1}}}, unas treinta veces más lenta que
la primera etapa de \omterm{def:b2:pt-paths:exhumation}{exhumación} de las rocas alpinas más profundas, y aproximadamente el doble de la velocidad de erosión
del modelo de engrosamiento.
\end{solution}
